\section{The complete numerical passage from the original tied family to a target window}
\label{sec:tied-window-completion}

\status{ZW1--ZW4 close the numerical and source-identification step between
the literal \texttt{ChainTied} supplied by the global theorem and WB's
positive character window.  The proof includes the weighted slope
calculation and every error absorption; these are not inferred merely
from the name of a source theorem.  The eleven arithmetic ledger
conditions remain fields of the supplied \texttt{ChainTied}.  No new
arithmetic existence assertion is added after that record.}

There are two different meanings of ``gap'' in a dependency discussion.
Producing the simultaneous geometric estimates is substantial arithmetic.
Once those estimates occur as the eleven fields of the literal tied
ledger, turning them into the asked-minus-main inequality is a finite
calculation.  We give that calculation here.  A single uniform constant $200R^3$ simplifies the error absorption
without changing the source's final exponent.

% BEGIN ZW1
\paragraph{ZW1: unpacking the original family without changing the host.}
\label{claim:ZW1}
Fix $\mathcal F:\texttt{FamP}(\epsilon)$ and suppose
\[
 \texttt{ChainTied}(\mathcal F,e_D,e_L,K,c_{HH},e_K,C_L,x).
 \tag{ZW.1}
\]
Its witnesses include one original index $i_0$, the frame
$\phi=(\mathcal F.D(x,i_0)).R$, a map
$\mathrm{pt}:\phi.A\to\mathcal F.I_x$, and the ledger $D$ with word
sets $W_u=(\mathcal F.D(x,\mathrm{pt}(u))).\mathrm{Word}$.
They also include the top set $T$, second set $T_2$, coefficients
$a_t,w_{u,j}$, finite cutoff $N$, actual logarithmic derivatives
$\ell$, and all of the following on these same objects:
\[
 \begin{gathered}
 D.x=x,\quad D.\epsilon=\epsilon,\quad
 C.e_{\rm Div}\le e_D,\quad C.e_{\rm Loss}\le e_L,
 \quad C.c_{HH}\le c_{HH},\\
 \texttt{Ident},\quad \texttt{SizeData}_2,\quad
 \texttt{ErrWindows},\quad\texttt{BNBounds},\\
 \texttt{GrowthInequalitySizes},\quad\texttt{ThinScales},\quad
 \texttt{CostWindows},\quad\texttt{RestSize}.
 \end{gathered}                                                    \tag{ZW.2}
\]
Here \texttt{SizeData}$_2$ means the eleven predicates in
\texttt{Formalism.LFunctionChainParameters.LChain}, not the different
same-named datum in local-ring cardinality.  Its eleven fields are
\[
 \begin{gathered}
 \texttt{HidaSetting},\quad \texttt{DIV},\quad
 \texttt{EComp5},\quad \texttt{HComp},\\
 \texttt{Count5},\quad \texttt{Slopes},\quad
 \texttt{QuotientSlopeBoundsWith},\\
 \texttt{LossBounds},\quad \texttt{SlopeIneq},\quad
 \texttt{HHBound},\quad \texttt{FParams}.
 \end{gathered}
\]
These are premises of the calculation below, not additional consequences
asserted here.

The source \texttt{TieArith} identifies membership in $T$ with membership
of the integer label in $\mathcal F.TTn(x)$.  The family field
\texttt{TT\_eq} gives exactly the same characterization for the top set
of the original host $H=(\mathcal F.D(x,i_0)).\mathrm{host}$.
Injectivity of the frame's integer labels therefore gives
$T=H.TT$, as sets rather than just equal cardinalities.  Its other
clause says that every element of $T_2$ is a digit of an original
second-set entry; such digits lie in $\phi.A_0$.  Finally
$\texttt{arithOf}(H)=\texttt{arith}(H)$ by their definitions.

Take the translate set to be $\phi.A$ and its coverage map to be the
identity.  It covers $\phi.A_0$.  Reindexing \texttt{Ident} by this map
gives precisely \texttt{IdentU}; \texttt{BNBounds} becomes
\texttt{BNBoundsU}; restricting the node bounds to $\phi.A_0$ gives
\texttt{SizesA}.  The word sets are finite and nonempty because their
original distinguished words are already fields of the datum.
\texttt{TieArith}'s last clause is the exact
\texttt{HostAddCharCompatible} equality at every primitive character,
including its linear centering term.  Consequently (ZW.1) supplies
\texttt{CoreAtH} on this original $H$ and its original $TT$.
The fields \texttt{TieGlobal} and \texttt{TieLocal} are not replaced by
new witnesses; the already supplied ledger is simply retained.
% END ZW1

% BEGIN ZW2
\paragraph{ZW2: the finite weighted slope calculation.}
\label{claim:ZW2}
Fix a translate and suppress it from notation.  The finite word set $W$
is nonempty, its weights $d_w$ are positive, and $\langle Z\rangle$
denotes $\sum_wd_wZ_w/\sum_wd_w$.  Put
\[
 A=\sum_ta_t,\quad n_T=\#T,\quad
 s=\langle\mu_{N,w}-h_w\rangle,\quad
 u=\max_w(\mu_{N,w}^{\max}-h_w)-s.
\]
Write $G=\mathrm{Ga}$, $F=\mathrm{Fo}$, $E_w=\mathrm{Ee}_w$,
$\mu_C=\mathrm{muCC}$, and $t_w=\log T_w$.  The inequality
$\mu_{N,w}\le\mu_{N,w}^{\max}$ shows $u\ge0$ by taking a
positive weighted average.

All coefficients below are the original components of $C$.  Define
\[
 \begin{aligned}
 B_1&=c_Sn_TaL_a(1+m/k),& B_2&=c_Sn_TaL_a,\\
 P&=c_\ell mB_0,&
 L_P&=c_Pn_Tmb\,x^{e_{\rm Loss}\epsilon},\\
 E_D&=c_Dn_TaL_a+c_Dm(B_0+N_c)x^{e_{\rm Div}\epsilon},\\
 V&=2mc_EB_0+c_H(n_TaL_a+mB_0)+Ac_C\log k,\\
 L_w&=c_\ell\mathrm{IO}
       +c_\ell n_Tm(a/k+b')t_w+c_\ell X''+1,\\
 M_L&=c_\ell L_P+c_\ell n_Tm(a/k+b')(c_CmL)+c_\ell X''+1,\\
 K_3&=B_1+B_2+V+P,\qquad U=K_3+M_L.
 \end{aligned}                                                     \tag{ZW.3}
\]
The subscripts $D,E,H,C,S,\ell,P$ abbreviate, in order,
\texttt{cDiv,cE,cH,cCount,cSlope,cSl1,cLoss}; $L_a=D.\texttt{La}$.
Thus $L_w,M_L,U$ are exactly the source's \texttt{Lo5},
\texttt{MLo5}, and \texttt{UBar5}.

The quotient and ordinary slope fields give
\[
 \begin{split}
 \mu_{N,w}&\ge-B_1+G+F-E_w-L_w,\\
 \mu_{N,w}^{\max}&\le B_2+G+F-E_w+P.
 \end{split}                                                       \tag{ZW.4}
\]
The exponent and height comparisons, followed by the variation clause of
\texttt{Count5}, give for any $v,w$
\[
 (-E_v-h_v)-(-E_w-h_w)\le V.                                      \tag{ZW.5}
\]
Indeed the exponent difference costs $2mc_EB_0$; the height difference
costs $c_H(n_TaL_a+mB_0)$ plus $A(t_v-t_w)$; and the last term is
at most $Ac_C\log k$.  Subtract the first inequality of (ZW.4) for
$w$ from the second for $v$, and average in $w$.  This gives
$\mu_{N,v}^{\max}-h_v-s\le K_3+\langle L_w\rangle$.
Taking the maximum in $v$ and using
$\mathrm{IO}\le L_P$, $t_w\le c_CmL$, and nonnegative coefficients,
we obtain
\[
             0\le u\le K_3+\langle L_w\rangle\le U.              \tag{ZW.6}
\]

Let
\[
 \Lambda=\sum_{j=1}^m\frac1f\sum_{\rho=1}^f
                 \log(w_{j\rho}v/\pi)+m\,\mathrm{lg}
                       -c_C L_a(1+m^2/k).
\]
The count field gives $t_w\ge\Lambda$.
Writing $S_D=\sum_t m\{\operatorname{ph}(a_t)-a_t\log\Omega\}$,
the divisibility inequality gives
$F-h_w\ge\mu_C-S_D+At_w-E_D$.
Combine this with (ZW.4), $A\ge0$, and
$E_w\le m(A\,\mathrm{lg}+c_EB_0)$, then average:
\[
 s\ge G+\mu_C-K_5-\langle L_w\rangle,\quad
 K_5=B_1+m(A\,\mathrm{lg}+c_EB_0)+S_D-A\Lambda+E_D.
\]
Expanding $\operatorname{ph}(y)=y\log y-y$, with the continuous
convention $0\log0=0$, and
$\log\Omega=\gamma+\log(\pi/v)-\mathrm{LpL}$, and distributing
the finite sums, gives exactly
\[
 S_D-A\Lambda+mA\,\mathrm{lg}
       =\mathrm{Main}+Ac_C L_a(1+m^2/k).
\]
Thus, with
\[
              K_0=Ac_C L_a(1+m^2/k)+E_D+mc_EB_0+B_1,
\]
we have $G\le\mathrm{Main}+K_0+\langle L_w\rangle+s-\mu_C$.

It remains to insert the original slope inequality.  Its hypotheses are
supplied by \texttt{FParams}, \texttt{HidaSetting} and
\texttt{HHBound}.  Write $\alpha=C.\texttt{cHH}$ for the ledger's
height exponent, distinct from its fixed global upper bound $c_{HH}$.
Put $H=\mathrm{HH}>0$ and $X=x^\alpha$.
For $u\ge0$, $H\le Xb$ and $b>0$,
\[
                u\log(1+H/u)\le(U+b)\log(1+X).                  \tag{ZW.7}
\]
For $u>0$, concavity of the logarithm with weights
$u/(u+b)$ and $b/(u+b)$ first bounds the left side by
$(u+b)\log(1+H/(u+b))$; use $H\le Xb$ and $u\le U$.
At $u=0$ the expression has value zero by continuous extension,
which also agrees with the source's totalized expression after
multiplication by $u$.

The parameters obey $0\le\delta\le\bar\delta<1$,
$0\le\eta\le\bar\eta$ and $0\le e\le\bar e\le1/4$.
The error terms in the slope inequality are nonnegative.  Hence we may
increase its reciprocal to $(1-\bar\delta)^{-1}$, replace its other
parameters by their upper bounds, and apply (ZW.7).
Finally $\mathrm{ASK}-G\le L_P$.  We have proved the full pointwise
estimate
\[
\begin{split}
 \mathrm{ASK}-\mathrm{Main}\le{}&K_0+\log|W|+L_P+M_L\\
 &+\frac{4\log(1+x^\alpha)(U+b)
       +(\bar\eta+4\sqrt{\bar e})x^\alpha b}
                  {1-\bar\delta}.
\end{split}                                                       \tag{ZW.8}
\]
The right side is exactly \texttt{ErrX5at}.  This proves its numerical
implication from all eleven original ledger fields, for every translate
separately, without a new mean or a choice of a better word.
% END ZW2

% BEGIN ZW3
\paragraph{ZW3: uniform absorption of every error.}
\label{claim:ZW3}
Assume the original \texttt{ErrWindows} and (ZW.2).
Fix $\epsilon>0$ and the global bounds $K,c_{HH},e_K,e_D,e_L$ before
the data.  Put
\[
 R=1+|K|+|c_{HH}|+|e_K|,\quad z=x^\epsilon,\quad
 E=n_Ta(1+m^2/k)+m(B_0+N_c).
\]
For all sufficiently large $x$, uniformly in every such ledger and
translate,
\[
                 \mathrm{ErrX5at}\le
                     x^{(e_D+e_L+3)\epsilon}E.                   \tag{ZW.9}
\]
We need only enlarge the threshold until
$x\ge3$, $\log x\le z$, and $200R^3\le z$.
These are eventual inequalities for each fixed positive $\epsilon$.

Here are the complete bounds.  \texttt{ErrWindows} says
$0<\epsilon\le1$, $n_T\ge1$, $b\ge1$, $xb\le a$,
$xm\le k\le a$, $1\le k\le x^{e_K}$, $kb'\le a$, $X''\le a$,
$0\le L,L_a\le2z$, $B_0,N_c\ge0$, and $0\le a_t\le2a$.
It also gives $m\le x^{1/2+\epsilon/2}\le x$ and
$\log|W|\le(k-3)\log2\le a$.
In particular
\[
 E\ge n_Ta\ge a\ge b\ge1,\quad E\ge mB_0,\quad
 A\le2n_Ta,\quad m/k\le1,\quad mb\le a.                         \tag{ZW.10}
\]
All seven multiplicative constants lie in $[0,K]\subset[0,R]$;
$0\le e_{\rm Div}\le e_D$ and $0\le e_{\rm Loss}\le e_L$.
The following table follows by direct substitution into (ZW.3):
\[
\begin{array}{c|l}
\text{quantity}&\text{upper bound}\\ \hline
 L_P&R E z^{e_L}\\
 V&7R^2 E z\\
 K_3&14R^2 E z\\
 M_L&7R^2 E z^{e_L+1}\\
 U&21R^2 E z^{e_L+1}\\
 K_0&11R E z+R E z^{e_D}.
\end{array}                                                       \tag{ZW.11}
\]
For clarity, the only potentially large mixed term is bounded by
\[
 c_\ell n_Tm(a/k+b')(c_CmL)
              \le4R^2 n_Ta(m^2/k)z\le4R^2Ez.
\]
Also $\log k\le e_K\log x\le Rz$, which bounds the final term
of $V$ by $2R^2Ez$.  The other two terms of $V$ cost at most
$5REz$.  The bounds for $K_3$ and $M_L$ then follow by addition;
$c_\ell X''+1\le(R+1)E\le2RE$.
These calculations prove every entry in (ZW.11), without a constant
depending on the number of words or on a changing field.

The remaining original error-window clause is
\[
 (\bar\eta+4\sqrt{\bar e})x^{C.c_{HH}}b
                         \le Kz\,n_Ta\le REz.                  \tag{ZW.12}
\]
Moreover $\bar\delta\le1/2$, so its reciprocal is at most $2$,
and, using $\alpha=C.\texttt{cHH}\le c_{HH}$,
$\log(1+x^\alpha)\le\log2+c_{HH}\log x\le Rz$.
Using (ZW.11)--(ZW.12), the last line of (ZW.8) is at most
\[
 8Rz\{21R^2Ez^{e_L+1}+E\}+2REz
                         \le178R^3Ez^{e_L+2}.
\]
Its first line is at most
$11REz+REz^{e_D}+E+REz^{e_L}+7R^2Ez^{e_L+1}$.
Because $z\ge1$ and $e_D,e_L\ge0$, the sum of both lines is
at most
\[
 200R^3 E z^{e_D+e_L+2}\le E z^{e_D+e_L+3},
\]
as required.  If the fixed bounds admit no ledger, the implication is
vacuous; otherwise their nonnegativity assertions follow from the
displayed original fields.  The proof is uniform in the original
indices, word sets, coefficients and host.
% END ZW3

% BEGIN ZW4
\paragraph{ZW4: the literal tied-family terminal theorem.}
\label{claim:ZW4}
Fix a cyclic target level $\Lambda=\mathrm{Level.cyc}(\psi)$,
$\epsilon>0$, and one family $\mathcal F:\texttt{FamP}(\epsilon)$.
Suppose $\texttt{EFamC}(\Lambda,\mathcal F.\texttt{toFamCore})$ holds
and there are fixed $e_D,e_L,K,c_{HH},e_K,C_L,x_0$ for which (ZW.1)
holds whenever $x\ge x_0$ and $x\in\mathcal F.\mathrm{Win}$.
Then, on all sufficiently late members of that family's windows,
\[
                    H_\psi(x)\le
                         x^{1/2+(e_D+e_L+7)\epsilon}.             \tag{ZW.13}
\]
Consequently the literal \texttt{FamilyClassConnected} conclusion for
$(\mathrm{GChi}\psi).\mathrm{cls}(3/2)$ suffices for the target's
zero-free half-plane.  No separate arithmetic-gap existence theorem is
needed after that conclusion.

Indeed ZW1 gives the exact \texttt{CoreAtH} ledger.  Its
\texttt{IdentU} field and the top/place bijections give the two
identities (WB.26), with the original $a_t,w_{u,j},\ell$ and $N$.
ZW2--ZW3 therefore bound the gap at every translate in the covered
subgroup.  The two \texttt{BNBoundsU} inequalities increase $E$ to
the bracket $B$ of (WB.22), term by term with nonnegative
coefficients.  \texttt{RestSize} supplies its lower size bounds;
the two upper bounds $m\le x^{1/2+\epsilon/2}$ and $xm\le k$
come from \texttt{ErrWindows}.  The remaining scalar upper bounds
$a\le2xb$ and $k\le n^2x^{1/2-\epsilon}$ are also fields of
\texttt{RestSize}, while $\#T_2/n\le1$ follows from $T_2\subset\phi.A$.
Its positive lower bounds imply $m\ge1$ and $\#T_2>0$, and
\texttt{ErrWindows} gives $n_T\ge1$.  Thus every scalar hypothesis
of WB6 holds, with the fixed gap exponent $C=e_D+e_L+3$.

The original \texttt{HostAx} supplies WB1--WB2's reciprocity and
ramification facts.  Its \texttt{EAtH} supplies the same primitive
retained characters, biased top-set counts, omitted-character ramp
bound, and nonretained logarithmic-derivative bound.  The finite sum
of the retained $|L'/L(1,\chi)|$ depends only on $\Lambda$, so gives
the uniform $K_0$ in WB5.  Enlarge the exceptional-character
constant to its maximum with zero, which preserves its bound for
$x\ge3$ and supplies WB5's nonnegativity convention.
The full \texttt{SizesA},
\texttt{ThinScales}, \texttt{CostWindows} and $T_2\subset A_0$
are exactly the remaining WB4 premises.  WB6 consequently yields
(ZW.13), with its additional threshold $x^\epsilon\ge1+4o$.

The definition of \texttt{EFamC} gives one host threshold and one
exceptional-character constant before $x$ and the family index.
The tied-family condition gives one ledger threshold and constants
before $x$.  Join these thresholds with ZW3's and WB's by a finite
maximum.  At each $x$, use the index $i_0(x)$ supplied by ZW1
for that $x$; no single index or host is fixed across the window.
The estimate holds at every real $x$ in the late part of each family
window, because every threshold precedes this index choice.

Finally fix any proposed zero with real part $\beta>1/2$ and choose
$0<\delta<\min(1/2,\beta-1/2)$.
\texttt{FamilyClassConnected} first fixes $e_D,e_L$ and an upper
range for $\epsilon$.  Choose a permitted $\epsilon>0$ so small
that $(e_D+e_L+7)\epsilon\le\delta$, and only then choose its
family and a window beyond the finite common threshold.  Since the
windows have fixed stretch $3/2>1$ and occur arbitrarily late,
WB8/AN3 contradict the proposed zero.  With
$e_D=2,e_L=0$ on the selected \texttt{topDesWin} route, the
pointwise exponent is exactly $1/2+9\epsilon$.

The arithmetic work retained here is precisely the production of the
literal common family and its links.  The numerical slope argument is
now complete even though its geometric slope premise
\texttt{SlopeIneq} still has to be installed by the upstream global
producer.  This distinction is reflected by separate implication and
producer statuses in the dependency audit.
% END ZW4
